\chapter{Récapitulatif}

\section{Les quarante-quatre briques, d'un coup d'œil}

\begin{center}
\scriptsize
\begin{tabularx}{\textwidth}{@{}clXX@{}}
\toprule
\textbf{\#} & \textbf{Outil} & \textbf{Son rôle ici} & \textbf{Écarté au profit de lui} \\
\midrule
\multicolumn{4}{@{}l}{\textbf{\color{insared}Couche 0 — Le socle}} \\
1 & Git & historique, branches, base de la CI & copier le dossier \\
2 & Python 3.10 & tout le code, version de la cible & Node, Go (agent à déployer) \\
3 & venv + épinglage & versions identiques partout & installation globale \\
\addlinespace
\multicolumn{4}{@{}l}{\textbf{\color{insared}Couche 1 — L'application web}} \\
4 & Flask & routes HTML et JSON, sessions & Django (trop de boîtes noires) \\
5 & Jinja2 & gabarits, échappement automatique & concaténation de chaînes \\
6 & Werkzeug & hachage pbkdf2, ProxyFix & hachage écrit à la main \\
7 & Gunicorn & 4 ouvriers, production & serveur de dév. de Flask \\
8 & nginx & TLS, statiques, en-têtes & Apache, Caddy \\
9 & Tailwind précompilé & apparence sans CDN & Tailwind par CDN (interdit par la CSP) \\
\addlinespace
\multicolumn{4}{@{}l}{\textbf{\color{insared}Couche 2 — Les données}} \\
10 & SQLAlchemy & 7 tables, pas d'injection SQL & SQL à la main \\
11 & Alembic & migrations versionnées & \cmd{ALTER TABLE} manuels \\
12 & SQLite & développement et tests & un serveur à lancer \\
13 & PostgreSQL & production, index partiels & MySQL \\
14 & Fernet & coffre des mots de passe root & chiffrement maison \\
\addlinespace
\multicolumn{4}{@{}l}{\textbf{\color{insared}Couche 3 — Les machines louées}} \\
15 & Docker & machines en moins d'une seconde & machines virtuelles \\
16 & Dockerfile & 6 images reproductibles & image figée à la main \\
17 & supervisord & 5 services, PID 1 correct & script shell, systemd \\
18 & OpenSSH & l'accès imposé par le sujet & --- \\
19 & ttyd & terminal sans rien installer & Shellinabox, Wetty \\
20 & Xvnc + XFCE + noVNC & bureau dans le navigateur & GNOME, KDE (trop lourds) \\
21 & entrypoint.sh & assemble la machine selon le mode & six images distinctes \\
\addlinespace
\multicolumn{4}{@{}l}{\textbf{\color{insared}Couche 4 — Le pilotage}} \\
22 & subprocess & la commande imposée, visible & SDK docker-py \\
23 & étiquettes Docker & identité des machines louées & préfixe de nom \\
24 & \cmd{docker -H ssh://} & piloter sans ouvrir de port & API Docker exposée \\
25 & l'agent & inscription et battement de cœur & liste de workers en dur \\
26 & l'ordonnanceur & nœud le moins chargé, déterministe & tirage aléatoire \\
27 & le Faucheur & expiration, réconciliation, reprise & fil dans l'application \\
28 & systemd + watchdog & services relancés, auto-réparation & cron, lancement manuel \\
\addlinespace
\multicolumn{4}{@{}l}{\textbf{\color{insared}Couche 5 — L'infrastructure}} \\
29 & Vagrant & 4 VM en une commande & création à la main \\
30 & VMware et 3 autres & fonctionne sur Intel et Apple & VirtualBox seul \\
31 & Ansible & 9 rôles idempotents & scripts shell \\
32 & Ansible Vault & secrets versionnables & secrets en clair, Vault externe \\
33 & inventaire dynamique & découvre les VM et leur clé & inventaire figé \\
\addlinespace
\multicolumn{4}{@{}l}{\textbf{\color{insared}Couche 6 — La qualité}} \\
34 & pytest & 168 tests en 20 s, sans Docker & unittest \\
35 & ruff & style et erreurs, en moins d'1 s & flake8 + pylint + isort \\
36 & bandit & motifs dangereux (SAST) & relecture humaine \\
37 & pip-audit & CVE des dépendances (SCA) & safety (base payante) \\
38 & gitleaks & secrets, historique compris & recherche manuelle \\
39 & nuclei + dast.sh & l'application en marche (DAST) & analyse statique seule \\
40 & hadolint + ansible-lint & l'infrastructure relue & relecture humaine \\
41 & Trivy & CVE des images, verrou & publication sans scan \\
42 & Syft & SBOM conservé & aucun inventaire \\
43 & GitHub Actions & 11 travaux, barrière & GitLab CI, Jenkins \\
44 & ghcr.io + compose & livrer, et lancer en une commande & registre tiers \\
\bottomrule
\end{tabularx}
\end{center}

\section{Qui réalise quoi dans le cahier des charges}

\begin{center}
\scriptsize
\begin{tabularx}{\textwidth}{@{}clX@{}}
\toprule
\textbf{\#} & \textbf{Fonctionnalité} & \textbf{Briques mobilisées} \\
\midrule
1 & Préparation, environnement & Git · Python · venv \\
2 & API Web Flask & Flask · Jinja2 · Werkzeug · Gunicorn · nginx · Tailwind \\
3 & Base de données & SQLAlchemy · Alembic · SQLite · PostgreSQL \\
4 & Conteneurisation Docker & Docker · Dockerfile · supervisord · entrypoint \\
5 & Infrastructure Vagrant & Vagrant · les quatre providers \\
6 & Automatisation Ansible & Ansible · Vault · inventaire dynamique \\
7 & Worker Agent & l'agent · les étiquettes · l'API des nœuds \\
8 & Resource Manager & l'ordonnanceur · l'index partiel de PostgreSQL \\
9 & Service de location & Flask · services.py · SQLAlchemy \\
10 & Déploiement des instances & subprocess · Docker · \cmd{docker -H ssh://} \\
11 & Accès SSH & OpenSSH · ttyd · Xvnc + XFCE + noVNC \\
12 & Gestion de l'expiration & le Faucheur · l'index \cmd{(status, end\_time)} \\
13 & Haute disponibilité & supervisord · \cmd{restart=always} · live-restore · l'agent · le Faucheur \\
14 & Sécurité DevSecOps & Fernet · bandit · pip-audit · gitleaks · nuclei · Trivy · Vault \\
15 & CI/CD & GitHub Actions · ghcr.io · Syft \\
16 & Tests d'intégration & pytest · docker compose · integration.sh \\
17 & Documentation & les schémas engendrés · ce document · le guide de 93 pages \\
\bottomrule
\end{tabularx}
\end{center}

\begin{retenir}
Aucune fonctionnalité n'est portée par un seul outil. La haute disponibilité,
par exemple, naît de cinq briques posées dans trois couches différentes :
supervisord dans la machine, \cmd{--restart=always} et \cmd{live-restore} sur le
nœud, l'agent et le Faucheur dans le parc. C'est pour cela qu'on ne peut pas
« ajouter la résilience à la fin ».
\end{retenir}

\section{Les sept ajouts au-delà du sujet}

\begin{center}
\small
\begin{tabularx}{\textwidth}{@{}XX@{}}
\toprule
\textbf{Ajout} & \textbf{Briques qui le rendent possible} \\
\midrule
Parc distribué à trois workers, alors qu'un hôte suffisait &
l'agent · l'ordonnanceur · \cmd{docker -H ssh://} \\
\addlinespace
Terminal web, en plus du SSH demandé &
ttyd · supervisord \\
\addlinespace
Bureau graphique complet dans le navigateur &
Xvnc · XFCE · noVNC · supervisord \\
\addlinespace
Coffre à secrets réaffichable après re-authentification &
Fernet · les sessions Flask · Ansible Vault \\
\addlinespace
Watchdog d'auto-réparation du service &
systemd · un timer d'une minute \\
\addlinespace
Inventaire dynamique indépendant du provider &
Vagrant · un script Python de 140 lignes \\
\addlinespace
Sauvegarde du travail avant destruction &
Docker \cmd{exec} · \cmd{tar} en flux · une table de plus \\
\bottomrule
\end{tabularx}
\end{center}

\section{Ce qu'il faut retenir}

\begin{description}
\item[Un outil se comprend par ce qu'on a écarté.] Savoir qu'on utilise Gunicorn
  ne vaut rien ; savoir que le serveur de Flask traite une requête à la fois,
  si. Chaque fiche de ce document a une section « pourquoi celui-là » pour cette
  raison.

\item[Les couches ne sont pas interchangeables.] On ne peut pas conteneuriser
  avant d'avoir une application, ni automatiser le déploiement avant d'avoir
  quelque chose à déployer. L'ordre dans lequel les briques ont été posées n'est
  pas un récit : c'est une contrainte.

\item[Les outils de la couche 6 ne servent à rien s'ils ne bloquent pas.]
  Un contrôle qui affiche un avertissement que personne ne lit est un coût sans
  bénéfice. Ici, aucune image n'atteint le registre si l'un des onze travaux
  échoue --- et c'est arrivé à nos propres images.

\item[Les meilleures briques viennent d'une panne.] Les étiquettes Docker
  existent parce que le Faucheur s'est supprimé lui-même. Le témoin de vie
  existe parce que la sonde du conteneur ne pouvait pas passer. L'inventaire
  dynamique existe parce que l'inventaire figé codait le provider en dur. Un
  projet sans incident n'a pas de garde-fous.
\end{description}

\vfill
\begin{center}
\color{gristexte}\small
Dépôt : \url{https://github.com/aziz3r/insacloud}\\
Pour la construction du projet pas à pas, voir le guide technique de 93 pages.
\end{center}
